\section*{第 \ref{ch:b1:structures} 章——代数结构}

\begin{solution}{exo:b1:structures:1}
\emph{稳定性：} $x * y = 1 \iff x + y - xy = 1 \iff (1-x)(1-y) = 0$，
当 $x, y \neq 1$ 时不可能成立。事实上，关键恒等式为
\[
1 - x * y = (1 - x)(1 - y):
\]
\omterm{def:b1:logic:map}{映射} $\varphi(x) = 1 - x$ 将 $(E, *)$ 映为 $(\R^*, \times)$，使得
$\varphi(x * y) = \varphi(x)\varphi(y)$——一个\omterm{def:b1:logic:inj}{双射}\omterm{def:b1:structures:morphism}{同态}。所有
公理均可传递：结合律与交换律来自
$\times$ 的相应性质；单位元为 $\varphi^{-1}(1) = 0$（验证：$x *
0 = x$）；$x$ 的逆元为 $\varphi^{-1}\bigl((1-x)^{-1}\bigr) =
1 - \frac{1}{1-x} = \frac{x}{x - 1}$（且它 $\neq 1$）。因此 $(E, *)$
构成\omterm{def:b1:structures:group}{阿贝尔群}。
\end{solution}

\begin{solution}{exo:b1:structures:2}
\begin{enumerate}
  \item 是：正数之积仍为正数，单位元为 $1$，逆元为
        $\frac 1x$，结合律由 $\R^*$ 继承。
  \item 否：不满足稳定性（$1 + 1 = 2 \notin \{-1,0,1\}$）。
  \item 是：标准的例子。
  \item 否：不满足稳定性（奇数 $+$ 奇数 $=$ 偶数），且没有单位元（$0$ 是
        偶数）。
\end{enumerate}
\end{solution}

\begin{solution}{exo:b1:structures:3}
记 $\mathrm{id}$，对换 $\tau_{12}, \tau_{13},
\tau_{23}$（交换所指定的两点），以及轮换 $c =
(1\,2\,3)$（即 $1 \mapsto 2 \mapsto 3 \mapsto 1$），及 $c^2 =
(1\,3\,2)$。下面给出 $\sigma\rho$ 的表（行为 $\sigma$，列为 $\rho$，
先作用 $\rho$）：
\begin{center}
\renewcommand{\arraystretch}{1.2}
\begin{tabular}{c|cccccc}
$\sigma\backslash\rho$ & $\mathrm{id}$ & $c$ & $c^2$ & $\tau_{12}$ &
$\tau_{13}$ & $\tau_{23}$ \\
\hline
$\mathrm{id}$ & $\mathrm{id}$ & $c$ & $c^2$ & $\tau_{12}$ &
$\tau_{13}$ & $\tau_{23}$ \\
$c$ & $c$ & $c^2$ & $\mathrm{id}$ & $\tau_{13}$ & $\tau_{23}$ &
$\tau_{12}$ \\
$c^2$ & $c^2$ & $\mathrm{id}$ & $c$ & $\tau_{23}$ & $\tau_{12}$ &
$\tau_{13}$ \\
$\tau_{12}$ & $\tau_{12}$ & $\tau_{23}$ & $\tau_{13}$ & $\mathrm{id}$
& $c^2$ & $c$ \\
$\tau_{13}$ & $\tau_{13}$ & $\tau_{12}$ & $\tau_{23}$ & $c$ &
$\mathrm{id}$ & $c^2$ \\
$\tau_{23}$ & $\tau_{23}$ & $\tau_{13}$ & $\tau_{12}$ & $c^2$ & $c$ &
$\mathrm{id}$ \\
\end{tabular}
\end{center}
不可交换的一对：$\tau_{12}\tau_{13} = c^2$，而
$\tau_{13}\tau_{12} = c$。（验算其中一个元素：
$\tau_{12}\tau_{13}$ 将 $1 \xmapsto{\tau_{13}} 3
\xmapsto{\tau_{12}} 3$, $3 \mapsto 1 \mapsto 2$, $2 \mapsto 2 \mapsto
1$：即 $1 \mapsto 3 \mapsto 2 \mapsto 1$，正是轮换 $c^2 =
(1\,3\,2)$.)
\end{solution}

\begin{solution}{exo:b1:structures:4}
$H$：$1 \in H$；对 $z, w \in H$，$\abs{zw^{-1}} = \abs z / \abs w =
1$：判别准则成立。$\R_+^*$：同理，把 $\abs{xy^{-1}}$ 换为
正性。并集：$\iu \in H$ 且 $2 \in \R_+^*$，但 $2\iu$ 的
模为 $2 \neq 1$，且它不是正实数：$2\iu \notin H \cup
\R_+^*$，故并集不满足稳定性——不是\omterm{def:b1:structures:subgroup}{子群}（正如
\cref{exo:b1:structures:6}所预言，两者互不包含）。
\end{solution}

\begin{solution}{exo:b1:structures:5}
\omterm{def:b1:structures:morphism}{同态}：$\eu^{\iu(\theta + \varphi)} =
\eu^{\iu\theta}\eu^{\iu\varphi}$（见\cref{thm:b1:complex:funceq}）。
核：$\eu^{\iu\theta} = 1 \iff \theta \in 2\pi\Z$，故 $\ker f =
2\pi\Z \neq \{0\}$，故不是\omterm{def:b1:logic:inj}{单射}。像：每个单位复数都可写成
$\eu^{\iu\theta}$ 的形式，其中 $\theta$ 适当选取（由极坐标形式可得），因此
$\operatorname{im} f = \mathbb{U}$，即单位圆。
$f$ 在长为 $2\pi$ 的半开区间上的限制，例如
$\intco{0}{2\pi}$ 或 $\intoc{-\pi}{\pi}$，是\omterm{def:b1:logic:inj}{单射}（两个像相同的角
之差是 $2\pi$ 的整数倍，且区间内每个\omterm{def:b1:logic:equiv}{等价类}恰有一个代表元
）；任何长度更大的区间都不行，
因为其中含有两个距离为
$2\pi$ 的点。
\end{solution}

\begin{solution}{exo:b1:structures:6}
交集：$e \in H \cap K$，且当 $x, y \in H \cap K$ 时，
$xy^{-1}$ 同时属于 $H$ 与 $K$。并集：若 $H \subseteq K$，则并集即为
$K$，是\omterm{def:b1:structures:subgroup}{子群}（另一情形同理）。反之，假设两个包含关系都不成立，
任取 $h \in H \setminus K$ 与 $k \in K \setminus
H$，并设 $H \cup K$ 是\omterm{def:b1:structures:subgroup}{子群}，则 $hk \in H \cup K$。
若 $hk \in H$，则 $k = h^{-1}(hk) \in H$，矛盾。若 $hk \in
K$，则 $h = (hk)k^{-1} \in K$，矛盾。因此 $H \cup K$ 不是
\omterm{def:b1:structures:subgroup}{子群}。
\end{solution}

\begin{solution}{exo:b1:structures:7}
首先注意，$x^2 = e$ 意味着 $x^{-1} = x$ 对每个 $x$ 都成立。于是对
$x, y \in G$：
\[
xy = (xy)^{-1} = y^{-1} x^{-1} = yx ,
\]
由\cref{prop:b1:structures:rules}~(2)可知，$G$ 是\omterm{def:b1:structures:group}{阿贝尔群}。
\end{solution}

\begin{solution}{exo:b1:structures:8}
$\Z/18\Z$ 的可逆元，即与 $18 = 2 \times 3^2$ \omterm{cor:b1:arith:bezout}{互质}的剩余类：
$\overline 1, \overline 5, \overline 7, \overline{11}, \overline{13},
\overline{17}$。$\overline 5$ 的逆元：$5 \times 11 = 55 = 3\times
18 + 1$，故 $\overline 5^{-1} = \overline{11}$。

$\overline 5 x = \overline 7$：乘以 $\overline{11}$，得 $x =
\overline{77} = \overline 5$（因 $77 = 4\times 18 + 5$）。唯一
解。

$\overline 6 x = \overline 3$：方程 $6x \equiv 3 \pmod{18}$
即 $18 \mid 6x - 3$。但 $6x - 3 = 3(2x - 1)$ 为奇数，而 $18$
为偶数：偶数不能\omterm{def:b1:arith:divides}{整除}奇数。无解。

$\overline 6 x = \overline{12}$：$6x \equiv 12 \pmod{18} \iff x
\equiv 2 \pmod 3$，解为 $x \in \{\overline 2, \overline 5,
\overline 8, \overline{11}, \overline{14}, \overline{17}\}$——共
六个。
\end{solution}

\begin{solution}{exo:b1:structures:9}
$\Z[\sqrt 2]$ 含 $0$ 与 $1$，且对减法
与乘法封闭：
\[
(a + b\sqrt 2)(c + d\sqrt 2) = (ac + 2bd) + (ad + bc)\sqrt 2 ,
\]
故它是 $\R$ 的子环（交换律、结合律、
分配律均继承而来）。可逆元：$(1 + \sqrt 2)(-1 + \sqrt 2) =
2 - 1 = 1$，故 $1 + \sqrt 2$ 可逆，其逆元为 $\sqrt 2 - 1
\in \Z[\sqrt 2]$。其幂 $(1 + \sqrt 2)^n$ 严格
递增（底数 $> 1$），故两两不同，且每个都
是可逆元（$\bigl((1+\sqrt2)^n\bigr)^{-1} = (\sqrt 2 - 1)^n$）。可逆元全体构成的
群是无限的。
\end{solution}

\begin{solution}{exo:b1:structures:10}
$x + x = (x + x)^2 = x^2 + x^2 + x^2 + x^2 = 4x^2 = 4x$——于是 $2x =
4x$，从而 $2x = 0$，即 $x + x = 0$（每个元素都是自身的
加法逆元）。于是
\[
x + y = (x+y)^2 = x^2 + xy + yx + y^2 = x + xy + yx + y ,
\]
从而 $xy + yx = 0$，即 $xy = -yx = yx$（用到 $-z = z$）。故 $A$ 是
交换的。

例：在 $\mathcal{P}(E)$ 上定义 $A + B = (A \cup B) \setminus (A
\cap B)$（对称差），并令 $A \times B = A \cap B$。容易
验证：$(\mathcal{P}(E), +)$ 是\omterm{def:b1:structures:group}{阿贝尔群}，零元为
$\emptyset$，每个\omterm{def:b1:logic:sets}{集合}都是自身的逆元；$\cap$ 满足结合律、
交换律，单位元为 $E$；分配律 $A \cap (B + C) = (A
\cap B) + (A \cap C)$ 成立（一个元素属于左边当且仅当它
属于 $A$，且恰属于 $B, C$ 之一）。又 $A \cap A = A$：每个
元素都是幂等的，符合要求。
\end{solution}

\begin{solution}{exo:b1:structures:11}
写出 $n$ 与 $n+1$ 时的假设：
\[
(xy)^{n+1} = x^{n+1} y^{n+1}
\quad\text{且}\quad
(xy)^{n+1} = (xy)(xy)^n = xy\,x^n y^n .
\]
比较得 $x^{n+1} y^{n+1} = x\,y\,x^n\,y^n$；在左边消去 $x$，
在右边消去 $y^n$，得 $x^n y = y x^n$。同样的计算用于次数高一次的
情形（$n+1$ 与 $n+2$）可得 $x^{n+1} y = y x^{n+1}$。于是
\[
y\,x^{n+1} = x^{n+1} y = x\,(x^n y) = x\,y\,x^n ,
\]
于是在右端消去 $x^n$（原式为 $y x \cdot x^n = x y \cdot x^n$）：
$yx = xy$。故 $G$ 是\omterm{def:b1:structures:group}{阿贝尔群}。
\end{solution}

\begin{solution}{exo:b1:structures:12}
\begin{enumerate}
  \item 设 $f \colon \Z \to \Z$ 为可加\omterm{def:b1:logic:map}{映射}，$a = f(1)$。由
        归纳法得 $f(k) = ka$（$k \in \N$），且 $f(-k) = -f(k) =
        -ka$，故 $f$ 是乘以 $a$ 的\omterm{def:b1:logic:map}{映射}。反之，
        每个\omterm{def:b1:logic:map}{映射} $k \mapsto ak$ 都是\omterm{def:b1:structures:morphism}{同态}：
        $(\Z,+) \to (\Z,+)$ 的\omterm{def:b1:structures:morphism}{同态}恰为
        某个固定整数的乘法。
  \item 设 $f \colon \Q \to \Z$ 为\omterm{def:b1:structures:morphism}{同态}，$x \in \Q$ 且
        $n \in \N^*$。则
        \[
        f(x) = f\Bigl(\underbrace{\tfrac xn + \dots +
        \tfrac xn}_{n}\Bigr) = n\,f\Bigl(\frac xn\Bigr) ,
        \]
        故整数 $f(x)$ 可被每个 $n \geq 1$ \omterm{def:b1:arith:divides}{整除}。
        满足此条件的整数只有 $0$：$f \equiv 0$。
\end{enumerate}
\end{solution}

\begin{solution}{pb:b1:structures:1}
\textbf{1.} \omterm{def:b1:counting:objects}{置换}是 $\intint1n$ 上的\omterm{def:b1:logic:inj}{双射}，即
$n$ 个对象的一个 $n$ 元排列：共有 $n!$ 个
（见\cref{thm:b1:counting:counts}）。取 $\sigma = [2,3,1]$，$\tau =
[1,3,2]$，$\sigma\tau$ 将 $1 \mapsto 1 \mapsto 2$，$2 \mapsto 3
\mapsto 1$，$3 \mapsto 2 \mapsto 3$：$\sigma\tau = [2,1,3]$；而
$\tau\sigma$ 将 $1 \mapsto 2 \mapsto 3$，$2 \mapsto 3 \mapsto
2$，$3 \mapsto 1 \mapsto 1$：$\tau\sigma = [3,2,1] \neq
\sigma\tau$。

\textbf{2.} 设 $\gamma, \gamma'$ 的支集 $S, S'$ 互不相交。
对 $x \in S$，有 $\gamma'(x) = x$ 且 $\gamma(x) \in S$，故
$\gamma\gamma'(x) = \gamma(x) = \gamma'\gamma(x)$。对称地，
对 $x \in S'$ 同理；两边都使每个 $x \notin S \cup S'$ 保持不动。故
$\gamma\gamma' = \gamma'\gamma$。

\textbf{3.} 对 $i \in \intint1n$，$i, \sigma(i),
\sigma^2(i), \dots$ 均落在有限集中，于是 $\sigma^a(i) =
\sigma^b(i)$，其中 $a < b$；由\omterm{def:b1:logic:inj}{单射}性得 $\sigma^{b-a}(i)
= i$，即序列回到 $i$。称\emph{轨道}为 $i$ 的
\omterm{def:b1:logic:sets}{集合} $\{i, \sigma(i), \dots, \sigma^{k-1}(i)\}$，其中 $k \geq 1$
为满足 $\sigma^k(i) = i$ 的最小者。两个轨道若交于一点，
则二者重合（都是该点的 $\sigma$ 正向像集），
故轨道构成 $\intint1n$ 的\omterm{thm:b1:logic:partition}{划分}；$\sigma$ 作用于每个
大小为 $k \geq 2$ 的轨道上，相当于 $k$-轮换 $(i\ \sigma(i)\ \cdots\
\sigma^{k-1}(i))$，并使单点不动。这些互不相交的
轮换之积处处与 $\sigma$ 一致。唯一性：在任一
分解为互不相交轮换的式子中，经过 $i$ 的轮换必须
为 $(i\ \sigma(i)\ \cdots)$——各轮换被迫是
带诱导作用的轨道。

\textbf{4.} 沿各轨道追踪：$1 \to 4 \to 2 \to 1$，$3 \to 5
\to 3$，$6 \to 7 \to 8 \to 6$：
\[
\sigma = (1\ 4\ 2)(3\ 5)(6\ 7\ 8) .
\]
若 $\sigma = \gamma_1\cdots\gamma_r$ 为互不相交轮换之积，各轮换的长度分别为
$k_1, \dots, k_r$，由交换性（第 2 题）得
$\sigma^m = \gamma_1^m\cdots\gamma_r^m$，又因各轮换的
支集互不相交，$\sigma^m = \mathrm{id}$ 当且仅当每个 $\gamma_i^m =
\mathrm{id}$，即 $k_i \mid m$ 对所有 $i$ 成立（$k$-轮换的阶为
$k$：$\gamma^m$ 把 $a_1$ 映为 $a_{1 + (m \bmod k)}$）。满足条件的最小
$m$ 为 $\operatorname{lcm}(k_1, \dots, k_r)$。本题中
$\operatorname{lcm}(3, 2, 3) = 6$。

\textbf{5.} 对每一点作用右端的式子，先作用最右边的
因子。$a_1 \mapsto a_2$ 由 $(a_1\ a_2)$ 给出，随后各因子
均不动 $a_2$，总体为 $a_1 \mapsto a_2$。对 $2 \leq i < k$：
$a_i$ 在 $(a_1\ a_i)$ 把它送到 $a_1$ 之前不受影响，而紧接着的
因子 $(a_1\ a_{i+1})$ 把 $a_1$ 送到 $a_{i+1}$，此后
再无因子移动它，总体 $a_i \mapsto a_{i+1}$。最后，$a_k$
除最左边的因子外，其余因子都不动它，最左因子将其送到 $a_1$。
这恰好是该轮换。由于每个\omterm{def:b1:counting:objects}{置换}都是若干轮换之积（第 3 题），
故它也是若干对换之积。对第 4 题的
$\sigma$：
\[
\sigma = (1\ 2)(1\ 4)\;(3\ 5)\;(6\ 8)(6\ 7),
\]
共五个对换。

\textbf{6.} 对 $b - a$ 作归纳。当 $b = a + 1$ 时，恒等式
显然成立（此时仅一个因子，$1 = 2\cdot1 - 1$）。当 $b > a + 1$ 时，
直接验证 $(a\ b) = (a\ \ a{+}1)\,(a{+}1\ \ b)\,(a\ \ a{+}1)$：
右端将 $a \mapsto a{+}1 \mapsto b \mapsto b$，$b
\mapsto b \mapsto a{+}1 \mapsto a$，$a{+}1 \mapsto a \mapsto a
\mapsto a{+}1$，其余不动。由归纳假设，$(a{+}1\ \ b)$ 是
由 $2(b - a - 1) - 1$ 个相邻对换组成的回文乘积，
故 $(a\ b)$ 为 $2(b - a) - 1$ 个相邻对换之积，
其个数为奇数。

\textbf{7.} $N(\mathrm{id}) = 0$，$\varepsilon = +1$。对 $(i\ \
i{+}1)$，唯一的逆序对为 $(i, i+1)$：$N = 1$，
$\varepsilon = -1$。对 $[2, 3, 1]$，其逆序对为 $(1,
3)$（取值 $2 > 1$）与 $(2, 3)$（取值 $3 > 1$）：$N = 2$，
$\varepsilon = +1$。

\textbf{8.} $\sigma$ 与 $\sigma\tau$ 的取值序列仅相差
位置 $i$ 与 $i + 1$ 上的取值互换。对不涉及 $i, i+1$ 的位置对，
均无变化。对 $k < i$，
位置对 $(k, i)$ 与 $(k, i+1)$ 交换了逆序状态
（同样两个取值与 $\sigma(k)$ 比较，只是位置顺序相反），
故二者贡献之和不变；
对 $k > i + 1$ 同理。唯一剩下的位置对 $(i,
i+1)$ 的逆序状态翻转。因此 $N(\sigma\tau) = N(\sigma) \pm 1$。

\textbf{9.} 设 $\tau = (a\ b)$ 为任一对换，由
第 6 题，它是奇数个相邻对换之积，
故右乘 $\tau$ 使 $N$ 改变一个
奇数（由第 8 题反复应用可得）：$\varepsilon(\sigma
\tau) = -\varepsilon(\sigma)$。现设 $\sigma = \tau_1\cdots
\tau_p$（均为对换），从恒等\omterm{def:b1:counting:objects}{置换}出发，经 $p$
次右乘，得 $\varepsilon(\sigma) =
(-1)^p\varepsilon(\mathrm{id}) = (-1)^p$。由于
$\varepsilon(\sigma)$ 由逆序定义，独立于任何分解，
故 $p$ 的奇偶性是 $\sigma$ 的
不变量。\omterm{def:b1:structures:morphism}{同态}性：设 $\sigma$ 写成 $p$ 个对换之积，$\sigma'$ 写成
$q$ 个对换之积，则 $\sigma\sigma'$ 由 $p + q$ 个对换构成：
$\varepsilon(\sigma\sigma') = (-1)^{p+q} =
\varepsilon(\sigma)\varepsilon(\sigma')$.

\textbf{10.} $k$-轮换是 $k - 1$ 个对换之积
（第 5 题），故 $\varepsilon = (-1)^{k-1}$。对一般的 $\sigma$
其轨道大小分别为 $k_1, \dots, k_r$（$k_i \geq 2$），另有 $f$
个不动点，此时 $c(\sigma) = r + f$，$n = k_1 + \dots + k_r + f$，
于是
\[
\varepsilon(\sigma) = \prod_{i=1}^r (-1)^{k_i - 1}
= (-1)^{\sum_i k_i - r} = (-1)^{n - f - r} = (-1)^{n -
c(\sigma)} .
\]

\textbf{11.} $\mathfrak A_n = \ker\varepsilon$ 是\omterm{def:b1:structures:subgroup}{子群}，因其为
\omterm{def:b1:structures:morphism}{同态}的核（\cref{def:b1:structures:morphism}）。取定一个对换 $\tau_0$
（$n \geq 2$ 时存在）。\omterm{def:b1:logic:map}{映射} $\sigma
\mapsto \sigma\tau_0$ 是 $\mathfrak S_n$ 的\omterm{def:b1:logic:inj}{双射}（其逆为自身），
它把 $\mathfrak A_n$ 与奇\omterm{def:b1:counting:objects}{置换}的
\omterm{def:b1:logic:sets}{集合}互换（第 9 题）。两\omterm{def:b1:logic:sets}{集合}\omterm{thm:b1:logic:partition}{划分} $\mathfrak S_n$
且大小相等：$\abs{\mathfrak A_n} = \frac{n!}2$。

\textbf{12.} \emph{逆序}（$[4, 1, 5, 2, 3, 7, 8, 6]$ 中）：
由取值 $4$ 得 $1, 2, 3$ 共三个；由 $5$ 得 $2, 3$ 共两个；
由 $7$ 得 $6$ 一个；由 $8$ 得 $6$ 一个。$N = 7$，
$\varepsilon = -1$。\emph{轮换型}：$c = 3$ 个轨道，$n = 8$：
$\varepsilon = (-1)^{8-3} = -1$。\emph{对换个数}：
第 5 题中共有五个对换：$(-1)^5 = -1$。三者一致。

\textbf{13.} $(a\ b)(a\ c)$（先作用最右边者）：$a \mapsto c
\mapsto c$；$c \mapsto a \mapsto b$；$b \mapsto b \mapsto a$，
故得 $3$-轮换 $(a\ c\ b)$。又 $(a\ c\ b)(a\ c\ d)$：$a \mapsto c
\mapsto b$；$b \mapsto b \mapsto a$；$c \mapsto d \mapsto d$；$d
\mapsto a \mapsto c$：即 $(a\ b)(c\ d)$，正如所言。

\textbf{14.} 设 $\sigma \in \mathfrak A_n$，由第 9 题，
$\sigma = \tau_1\cdots\tau_{2m}$ 是偶数个
对换之积。将其按相邻两个分组：
$\tau_{2i-1}\tau_{2i}$：若二者相等，则此对
为恒等\omterm{def:b1:counting:objects}{置换}而消去；若二者恰有一个公共点，则
由第 13 题的第一个恒等式，将该对写作一个 $3$-轮换；
若二者不交，则由第二个恒等式将其写作两个
$3$-轮换。因此 $\sigma$ 是若干 $3$-轮换之积（或
恒等\omterm{def:b1:counting:objects}{置换}，即空积；对 $n \geq 3$ 还可取 $(1\ 2\
3)^3$）。

\textbf{15.} $(1\ 2)(3\ 4) = (1\ 3\ 2)(1\ 3\ 4)$（第 13 题，
取 $a{=}1, b{=}2, c{=}3, d{=}4$）。对 $5$-轮换，由
第 5 题，$(1\ 2\ 3\ 4\ 5) = (1\ 5)(1\ 4)(1\ 3)(1\ 2)$，并
两两配对：$(1\ 5)(1\ 4) = (1\ 4\ 5)$，$(1\ 3)(1\ 2) = (1\ 2\ 3)$：
\[
(1\ 2\ 3\ 4\ 5) = (1\ 4\ 5)(1\ 2\ 3) .
\]
（验算 $3$：$(1\ 2\ 3)$
将 $3 \to 1$，再经 $(1\ 4\ 5)$ 将 $1 \to 4$，总效果为 $3 \to 4$，正确。）

\textbf{16.} 两边同时作用于任一点。对 $i =
\sigma(a_j)$ 的情形：左边为 $\sigma\bigl((a_1\ \dots\
a_k)(a_j)\bigr) = \sigma(a_{j+1})$（下标按模 $k$ 理解），这正是
右边对 $\sigma(a_j)$ 的作用。对不具此形式的 $i$：
$\sigma^{-1}(i)$ 在支集之外，故左边
使 $i$ 不变，右边亦然。两边处处相等。

\textbf{17.} 把格子 $c'$ 中的方块滑入空格 $c$，
这交换了格子 $c$ 与 $c'$ 的内容（方块 $9$，
即空格，移至 $c'$）。若方块 $\sigma(i)$ 原位于格子 $i$，则
新位置为 $\sigma' = \sigma \circ (c\ c')$：各格子内容与原先相同，
唯格子 $c, c'$ 中的内容互换。由
第 9 题，$\varepsilon(\sigma') = -\varepsilon(\sigma)$。

\textbf{18.} 一次移动把空格移至相邻格子，其所在的行或列
恰改变 $1$，故空格的出租车距离
$d$（到格子 $9$）改变 $\pm1$，于是 $(-1)^d$ 变号。由于每次移动
都使 $\varepsilon(\sigma)$ 与 $(-1)^{d(\sigma)}$ 同时变号，
故其乘积 $I(\sigma)$ 经每次移动都保持不变：
是一个不变量。

\textbf{19.} 复原局面的 $\varepsilon = +1$，$d = 0$，
故 $I = +1$。目标局面（方块 $7, 8$ 互换，空格归位）恰为
格子 $7$ 与 $8$ 中内容的对换：$\varepsilon
= -1$，$d = 0$，故 $I = -1$。由于 $I$ 不变，两者
取值不同，故任何移动序列都无法将二者相连。

\textbf{20.} 空格在原位的局面，是
$8$ 个方块在格子 $1, \dots, 8$ 之间的一个\omterm{def:b1:counting:objects}{置换}，即
$\mathfrak S_8$ 中的元素；其 $d = 0$，故 $I = \varepsilon(\sigma)$。
可达要求 $I = +1$，即 $\sigma \in \mathfrak A_8$；
在本阶段不加证明地承认的逆命题断言，$\mathfrak A_8$ 的全部元素均可达。计数：
$\abs{\mathfrak A_8} = \frac{8!}2 = 20\,160$（第 11 题）。

\textbf{21.} 由第 20 题，空格在原位时的可达排列
恰构成 $\mathfrak A_8$——特别地，它是 $\mathfrak S_8$ 的一个
\emph{\omterm{def:b1:structures:subgroup}{子群}}：两个可解打乱的复合，
或取其逆，仍可解，这一点远非
仅凭纯粹的谜题推理便可看出。

\textbf{22.} (a) 空格在原位时，方块的一个 $3$-轮换：
$\varepsilon = +1$（第 10 题），$d = 0$，故 $I = +1$：可达
（由在本阶段不加证明地承认的逆命题）——即三个方块可以轮换。
(b) 方块 $5$ 与空格互换，此局面是
格子 $5$ 与 $9$ 中内容的对换，故
$\varepsilon = -1$；空格位于中心，与原位的出租车距离为
$d = 2$，故 $(-1)^d = +1$，$I = -1$：
不可达。不能简单地“把空格停在中间”，
而让其余方块仍然排好序。

\textbf{23.} 设 $f \colon \mathfrak S_n \to \{\pm1\}$ 为一个
\omterm{def:b1:structures:morphism}{同态}。对任意两个对换 $\tau, \tau'$，由第 16 题，
存在 $\sigma$ 使 $\sigma\tau\sigma^{-1} = \tau'$（将
两个被移动的点映到另外两点上；$n \geq 3$ 保证有足够的余地
以便做到这一点，尽管 $n = 2$ 时此处亦平凡）。于是 $f(\tau') =
f(\sigma)f(\tau)f(\sigma)^{-1} = f(\tau)$，因 $\{\pm1\}$ 是
阿贝尔（交换）的：$f$ 在对换上为常值。若该常值为
$+1$，则 $f = 1$ 于一切对换之积上，
即处处成立（第 5 题）。若其为 $-1$，则 $f(\sigma) =
(-1)^p = \varepsilon(\sigma)$，即在 $p$ 个对换之积上。
故 $f \in \{1, \varepsilon\}$。

\textbf{24.} (i) $\varepsilon$ 的\omterm{def:b1:structures:morphism}{同态}性质与
核的理论赋予 $\mathfrak A_n$ \omterm{def:b1:structures:subgroup}{子群}结构及其
大小，而与\cref{prop:b1:structures:kernel}相似的推理
贯穿第 11 题与第 21 题。(ii) 计数：
$\abs{\mathfrak S_n} = n!$，第 11 题的折半论证，
以及第 20 题的计数 $20\,160$ 均是
\cref{ch:b1:counting}中方法的应用。(iii) 第 8--9 题解决了一个
真正的良定义问题——“对换个数的
奇偶性”预设了这一奇偶性不依赖于分解，正如
$\Z/n\Z$ 的运算在\cref{def:b1:structures:zn}中要求
与代表元的选取无关。

\textbf{25.} 符号是一个取值于 $\{\pm1\}$ 的
单一计算，只需证明一次其为良定义，它同时发挥三重
作用：在内部，它将 $\mathfrak S_n$ 一分为二，并
分离出 $\mathfrak A_n$，并给出其 $3$-轮换生成元；
在外部，它可在一行之内判定一个问题（“这两个
方块能否互换？”），这是朴素搜索永远无法解决的，因为
任何有限个失败的移动序列都不足以证明不可能性；而
从结构上看，它是
行列式公式 $\det A = \sum_\sigma \varepsilon(\sigma)\,
a_{1\sigma(1)}\cdots a_{n\sigma(n)}$ 中正负交替的引擎，见\cref{ch:b1:det}。这一
模式——找出被每一步基本移动保持不变的量，
并在起点与目标处分别计算它——是
数学家应对“是否可能？”之类问题的
标准武器，每当群作用于一个状态\omterm{def:b1:logic:sets}{集合}时，它
便会再次出现。
\end{solution}
